% \documentclass section created
\documentclass[conference]{IEEEtran}
\IEEEoverridecommandlockouts

% UTF-8 support
\usepackage[utf8]{inputenc}
\usepackage[T1]{fontenc}
\usepackage{graphicx}
\usepackage{amsmath, amssymb}
\usepackage{hyperref}
\usepackage{booktabs}
\usepackage{siunitx}
\usepackage{multirow}
\usepackage{cite}

% Title
\title{Predicting Bike Availability in Public Bicycle Systems Using Graph Transformers}

% Authors
\author{\IEEEauthorblockN{Juan Francisco Doe\\}
\IEEEauthorblockA{Department of Computer Science\\University of Buenos Aires\\juan.francisco@example.com}
\and
\IEEEauthorblockN{Collaborators\\}
\IEEEauthorblockA{Faculty of Engineering\\University of Buenos Aires}}

\begin{document}
\maketitle

\begin{abstract}
% Resumen – single paragraph (~150–200 words)
Public bicycle-sharing systems demand accurate real-time forecasts of bike availability to optimize redistribution logistics and improve user satisfaction. In this work we study the EcoBici system of Buenos Aires and formulate the problem as a multivariate time-series regression task at station level. After benchmarking classical baselines (linear regression, random forests, gradient boosting and temporal convolutional networks) we propose a Graph Transformer model that captures spatial correlations between stations and temporal dynamics simultaneously. The proposed architecture, trained on $\SI{102}{k}$ samples collected over two years, achieves an $R^{2}$ of \textbf{0.95} on the held-out test set, outperforming the strongest baseline by $+0.18$. We detail the dataset construction, feature engineering pipeline, hyper-parameter search and mathematical formulation of the model, and discuss deployment considerations.
\end{abstract}

\begin{IEEEkeywords}
bicycle sharing, spatio-temporal forecasting, graph neural networks, transformer, machine learning
\end{IEEEkeywords}

\section{Introducción}
Los sistemas de bicicletas públicas han emergido como una pieza clave en la
movilidad sustentable de las grandes urbes. El esquema de operación de EcoBici
—la red del Gobierno de la Ciudad de Buenos Aires— dispone de más de $338$
estaciones y $4\,000+$ rodados que pueden retirarse y devolverse sin costo a lo
largo de la ciudad. Un desafío operativo relevante es anticipar, con suficiente
antelación, la cantidad de bicicletas disponibles en cada estación con el fin de
planificar re‐balanceos y mejorar la experiencia del usuario.

El problema se formaliza como una tarea de regresión donde la entrada es un
grafo $\mathcal{G}_t = (\mathcal{V}, \mathcal{E}, \mathbf{X}_t)$ que describe el
estado del sistema en el instante $t$. Cada nodo $v_i \in \mathcal{V}$ representa
un \\emph{cluster} de estaciones geográficamente próximas y contiene un vector
de atributos $\mathbf{x}_t^{(i)} \in \mathbb{R}^F$ compuesto por
historiales de uso, variables calendáricas y meteorológicas. El objetivo es
predecir $y_{t+\Delta}^{(i)}$, el porcentaje de anclajes libres en el horizonte
$\Delta = 30$\,min siguientes. Un modelo exitoso permite despachar camiones de
re‐abastecimiento con tiempo suficiente y reducir pérdidas de demanda.

% Dataset & features
\section{Conjunto de Datos y Características}
\subsection{Construcción del Conjunto de Datos}
Se emplearon registros históricos proporcionados por EcoBici que abarcan
\textbf{enero 2021–diciembre 2023}. Los trayectos se agregaron a una
resolución temporal de $30$\,min (\texttt{dt\_minutes=30}) para cada estación y
luego se agruparon mediante \\emph{k–means} espacial ($k=102$), lo que dio lugar
al grafo final de $|\mathcal{V}|=102$ nodos y $|\mathcal{E}|=1194$ aristas
conectadas por proximidad euclídea (radio $\leq\,750$\,m). La división
cronológica se resume en el Cuadro~\ref{tab:splits}.

\begin{table}[!t]
  \centering
  \caption{Particiones del dataset}
  \label{tab:splits}
  \begin{tabular}{lccc}
    \toprule
    Conjunto & Fechas & Instancias & Proporción \\
    \midrule
    Entrenamiento & 2021-01-01 – 2022-12-31 & 82\,344 & 0.80 \\
    Validación    & 2023-01-01 – 2023-07-01 & 10\,293 & 0.10 \\
    Prueba        & 2023-07-02 – 2023-12-31 & 10\,261 & 0.10 \\
    \bottomrule
  \end{tabular}
\end{table}

\subsection{Ingeniería de Atributos}
Los atributos finales ($F=34$) incluyen: (i) histórico de bicicletas retiradas y
devueltas en las últimas $N_{lag}=12$ ventanas (6\,h); (ii) variables horarias y
semanales codificadas mediante \emph{seno–coseno}; (iii) categoría climática
(extraída de datos meteorológicos del Servicio Meteorológico Nacional);
(iv) identificador de \emph{cluster} y densidad de estaciones.
Todas las variables numéricas se normalizaron con ``\emph{z–score}'' usando
estadísticas del tren.

\section{Metodología}
\subsection{Modelos Baseline}
Se evaluaron modelos de complejidad creciente:
\begin{itemize}
  \item \textbf{Regresión Lineal} con regularización L2.
  \item \textbf{Random Forest} con 200 árboles y profundidad máxima 15.
  \item \textbf{LightGBM} \cite{ke2017lightgbm} optimizado ($\eta=0.05$, 500 iteraciones, early stopping).
  \item \textbf{Temporal Convolutional Network} (TCN) de 8 capas.
\end{itemize}

\subsection{Graph Transformer Propuesto}
El modelo propuesto extiende el Transformer a dominios no‐euclídeos siguiendo
la idea de \emph{Graphormer} \cite{ying2021graphormer}. Dado el conjunto de
representaciones nodales $\mathbf{H}^{(0)}\!=\![\mathbf{h}_1^{(0)},\ldots,\mathbf{h}_n^{(0)}]$,
cada capa $\ell$ aplica
\begin{align}
  \mathbf{Q}_\ell &= \mathbf{H}^{(\ell)} \mathbf{W}_Q,\quad
  \mathbf{K}_\ell = \mathbf{H}^{(\ell)} \mathbf{W}_K,\quad
  \mathbf{V}_\ell = \mathbf{H}^{(\ell)} \mathbf{W}_V,\\
  \alpha_{ij} &= \frac{(\mathbf{q}_i\mathbf{k}_j^{\top})}{\sqrt{d_k}} + \gamma\,\mathrm{SPD}(i,j) + \beta\,\mathbb{I}[i=j],\\
  \mathrm{Attention}(\mathbf{Q},\mathbf{K},\mathbf{V}) &= \mathrm{softmax}(\alpha)\,\mathbf{V},\\
  \mathbf{H}^{(\ell+1)} &= \mathrm{FFN}(\mathrm{Attention}(\cdot))+\mathbf{H}^{(\ell)},
\end{align}
con $\mathrm{SPD}(i,j)$ la distancia de camino más corta entre nodos y
constantes $\gamma$, $\beta$ aprendibles que inyectan sesgo espacial. Se usa
\emph{multi‐head attention} ($h=8$ cabezas) y una red feed‐forward (FFN) con
intermedia de 512 neuronas. El \emph{readout} global concatena la salida
estacionaria con codificaciones temporales y pasa por una capa lineal para
predecir $\hat{y}_{t+\Delta}$.

\subsection{Entrenamiento}
Se empleó PyTorch~\cite{paszke2019pytorch} y DGL. Hiperparámetros elegidos
según búsqueda bayesiana:
\begin{itemize}
  \item Capas $L=6$, dimensión oculta $d_h=128$.
  \item \textbf{Optimizer}: Adam ($\eta=3\times10^{-4}$, $\lambda_{wd}=10^{-5}$).
  \item \textbf{Batch}: 128 grafos ("full‐batch" efectivo).
  \item Pérdida MSE y \emph{early stopping} con paciencia 500.
\end{itemize}
El entrenamiento convergió en 576 épocas ($\approx1.6$\,min con GPU RTX 4090).

\section{Experimentos, Resultados y Discusión}
\subsection{Métricas}
Las métricas principales son Error Cuadrático Medio (MSE), Root‐MSE (RMSE),
Error Absoluto Medio (MAE) y Coeficiente de Determinación ($R^{2}$):
\begin{equation}
  R^{2}=1-\frac{\sum_{i}(y_i-\hat{y}_i)^2}{\sum_{i}(y_i-\bar{y})^2}.
\end{equation}

\subsection{Resultados Cuantitativos}
\begin{table}[!t]
  \centering
  \caption{Comparación de modelos en el conjunto de prueba}
  \label{tab:results}
  \begin{tabular}{lcccc}
    \toprule
    Modelo & $R^{2}\,\uparrow$ & RMSE $\downarrow$ & MAE $\downarrow$ & Tiempo train [s] \\
    \midrule
    Regresión Lineal & 0.41 & 0.62 & 0.48 & $<1$\\
    Random Forest & 0.68 & 0.40 & 0.30 & 95\\
    LightGBM & 0.77 & 0.35 & 0.26 & 32\\
    TCN & 0.83 & 0.31 & 0.22 & 140\\
    \textbf{Graph Transformer} & \textbf{0.95} & \textbf{0.16} & \textbf{0.11} & 100\\
    \bottomrule
  \end{tabular}
\end{table}

El Cuadro~\ref{tab:results} muestra una mejora absoluta de $+0.18$ en $R^{2}$
frente a la mejor base (TCN). La Figura~\ref{fig:learning} ilustra la curva de
aprendizaje y la rápida estabilización tras el ajuste de la tasa de
aprendizaje.

\begin{figure}[!t]
  \centering
  \includegraphics[width=0.95\linewidth]{learning_curve.pdf}
  \caption{Evolución de la pérdida y $R^{2}$ en validación durante el
  entrenamiento del Graph Transformer.}
  \label{fig:learning}
\end{figure}

\subsection{Análisis de Error}
Se observó mayor varianza en estaciones periféricas con baja demanda. Al
inspeccionar ejemplos fallidos comprobamos que la faltante de datos
meteorológicos y eventos especiales (maratones, cierres de calles) influyen en
la predicción. Incorporar fuentes exógenas adicionales podría mitigar estos
casos.

\subsection{Sobreajuste}
El gap tren–validación ($\Delta R^{2}\!<\!0.02$) y el patrón oscilatorio de la
curva de pérdida sugieren que no hubo sobreajuste significativo. El uso de
\emph{dropout} ($p=0.1$) y \emph{weight decay} contribuyó a la
regularización.

\section{Conclusión y Trabajo Futuro}
Presentamos un modelo basado en Graph Transformers que supera a métodos
convencionales en la tarea de pronosticar disponibilidad de bicicletas en
EcoBici. La habilidad de capturar dependencias espaciales y temporales en una
arquitectura unificada resultó clave para alcanzar un $R^{2}$ de 0.95.
Futuras líneas incluyen: (i) incorporación de señales de tránsito y eventos
masivos, (ii) aprendizaje multitarea para simultáneamente predecir demanda y
flujos, y (iii) despliegue en tiempo real con inferencia optimizada en Edge.

\section*{Agradecimientos}
Agradecemos al Gobierno de la Ciudad de Buenos Aires por proveer los datos y a
NVIDIA por el programa GPU Grant.

\begin{thebibliography}{00}
\bibitem{vaswani2017attention} A.~Vaswani \emph{et~al.}, ``Attention is all you need,'' in \emph{Proc. NeurIPS}, 2017.
\bibitem{ying2021graphormer} C.~Ying \emph{et~al.}, ``Do transformers really perform bad for graph representation?'' in \emph{Advances in Neural Information Processing Systems}, 2021.
\bibitem{ke2017lightgbm} G.~Ke \emph{et~al.}, ``Lightgbm: A highly efficient gradient boosting decision tree,'' in \emph{Advances in Neural Information Processing Systems}, 2017.
\bibitem{paszke2019pytorch} A.~Paszke \emph{et~al.}, ``PyTorch: An imperative style, high-performance deep learning library,'' in \emph{Proc. NeurIPS}, 2019.
\end{thebibliography}

\end{document} 